% \vspace{-0.5cm}
\section{Conclusions and Future Work}
This paper presents \tool, a unified framework for DL test dataset adequacy assessment.
The framework integrates representative adequacy metrics in a consistent execution workflow while preserving their metric-specific processing requirements.
By providing extensible mechanisms for configuration, metric processing, and reporting, \tool simplifies the execution, comparison, and deployment of diverse adequacy techniques.
Additionally, by minimizing the engineering effort required to evaluate existing metrics, \tool enables researchers and practitioners to easily integrate test adequacy assessment into the development lifecycle of the DL system.

% As DL models continue to evolve through retraining, fine-tuning, and deployment optimization, practitioners are increasingly faced with large test datasets and limited labeling resources. In such settings, the assessment of the adequacy can provide valuable guidance for selecting, prioritizing, and evaluating test datasets. By reducing the engineering effort required to configure and execute existing adequacy metrics, \tool enables researchers and practitioners to more easily incorporate adequacy assessment into the DL development lifecycle.

% For future work, we plan to further improve the extensibility of the framework by introducing a lightweight plug-in mechanism that allows the newly proposed adequacy metrics to be integrated with minimal implementation effort.
As new adequacy techniques continue to emerge, we envision \tool as a continuously evolving platform that facilitates the reproducible evaluation and practical adoption of DL test adequacy assessment methods.
As future work, we plan to incorporate more metrics in the framework.
